\chapter{Quantum Mechanics}\label{quantum-mechanics}
\chapterstatus{planned}

In quantum mechanics states are rays in a Hilbert space, observables are self-adjoint operators and
symmetries are unitary representations. This chapter states the axioms precisely and solves the basic
examples with the spectral theory and representation theory developed earlier in the book.

\noindent\textbf{Prerequisites.} Chapter~\ref{functional-analysis} (Functional Analysis), Chapter~\ref{operator-algebras} (Operator Algebras), Chapter~\ref{harmonic-analysis} (Harmonic Analysis), Chapter~\ref{sobolev-spaces} (Partial Differential Equations), Chapter~\ref{representation-theory} (Representations of Groups and Lie Algebras), Chapter~\ref{lie-groups} (Lie Groups), Chapter~\ref{probability} (Probability Theory) and Chapter~\ref{classical-mechanics} (Classical Mechanics).

\noindent\textbf{Planned contents.}
\begin{enumerate}
\item The axioms: states, observables, the Born rule, expectation values; the Heisenberg uncertainty
principle; mixed states and density operators; Gleason's theorem (statement).
\item Dynamics: Stone's theorem, the Schr\"odinger equation, the Schr\"odinger and Heisenberg pictures;
the canonical commutation relations and the Stone--von Neumann theorem.
\item Self-adjointness of Schr\"odinger operators: the Kato--Rellich theorem; free particle, the
harmonic oscillator (Hermite functions, ladder operators), the particle in a box.
\item Angular momentum and $\mathrm{SU}(2)$: spin, addition of angular momenta, Clebsch--Gordan
coefficients; the hydrogen atom and its hidden $\mathrm{SO}(4)$ symmetry.
\item Spectral theory of Schr\"odinger operators: discrete and essential spectrum, bound states,
the min-max principle, scattering states (survey).
\item Approximation methods: perturbation theory (Rayleigh--Schr\"odinger), the variational method,
the WKB approximation.
\item Composite systems: tensor products, identical particles, bosons and fermions, Fock space;
entanglement and Bell's inequalities.
\item Quantisation: from Poisson brackets to commutators, the Groenewold--van Hove theorem, Weyl
quantisation (survey).
\item Supersymmetric quantum mechanics: graded Hilbert spaces (Section~\ref{multilinear-algebra:section-super}), supercharges $Q$ with $H = Q^2$, pairing of bosonic and fermionic states, the
Witten index $\operatorname{str}e^{-tH}$ and its invariance; Witten's deformation $d_t = e^{-tf}de^{tf}$ and the Morse inequalities.
\item Exercises with solutions.
\end{enumerate}
